\documentclass[11pt]{article}
\usepackage{amsmath,amssymb,amsthm,geometry}
\geometry{margin=.9in}\setlength{\emergencystretch}{1.5em}
\newtheorem{theorem}{Theorem}
\newcommand{\dd}{\,d}
\title{Core output descent and the pre-weak probe minimum}
\author{AN04 proof increment}\date{Version 1}
\begin{document}\maketitle

\paragraph{Domains and conclusions.}
Fix $c>0$, $\rho\in[13/20,7/10]$, $\lambda=\rho^2$,
$\omega=(1-\lambda)/2$, $L=\lambda/2$, and $\zeta>a$ where
$a=\rho K(\rho a)$, $K(t)=\phi(t)+t\Phi(t)$. Use the aligned balanced
isotropic Gaussian flow with $S_0=cq^{1/\omega}$. The origin
$s=t-t_h=0$ is the actual first right-ancestry half mass in normalized
time. All limiting statements below concern any fixed-horizon
diagonal extraction of this family as in BRW. Write
$D(s)=e^{Ls}/(1+s)$, $P_x=\Phi(\rho x)$, and $r_+=\max(r,0)$.

\begin{theorem}[Logistic resident limit]\label{an04codv1:main}
There are $C<\infty$, $k_y,k_E>0$ independent of the extracted law
such that, for all $s\ge0$,
\[
 y_a(s)\le C-k_yD(s),\qquad
 \mathcal E_\zeta(s)\le C-k_ED(s).
\]
The positive part of the active Q1 output exposure is $o(D(s))$,
uniformly over those laws.
\end{theorem}

\paragraph{Source inventory.}
HBD denotes the trained half-mass increment, prefix \texttt{an04hbdv1:};
BRW denotes the return increment, prefix \texttt{an04brwv1:}.
HBD \texttt{clock, incoming, outer, physical} concern the critical
family at fixed nominal mass $A=S_0e^t$, with compact Q1/Q4 labels
for individual chart assertions, and all labels for the Cartesian
bound. BRW extends extraction to each fixed finite resident horizon.
In the limit $N=0$, $M=(1+e^{-s})^{-1}$ and $\dd\mu=M\dd\pi$ with
normalized right label marginal $\dd\pi=4\cos^2\alpha\,\dd\alpha/(2\pi)$.
Q4 is one common bottom path $(x_a,y_a)$ with normalized mass $1/2$.
Q1 paths are co-ordered by their initial angle in this same field.
These claims follow from HBD entrance and common-field order as
continued in BRW; they do not require ordering two separate laws.
Put $Y=\int y\dd\mu$, $X_m=\int x\dd\mu$, $K_s=\int K(\rho x)\dd\mu$,
$S_x=\int F(\rho x,\rho\tilde x)\dd\mu$,
$F(t,r)=\phi(\min(t,r))+r\Phi(\min(t,r))$. BRW field is
\begin{align}
 2x'&=-(1-M)x+\rho\phi(\rho x)-\rho S_x+\lambda yP_x,\nonumber\\
 2y'&=-(1-M)y-Y+\rho K(\rho x),\qquad
 2Y'=(1-2M)Y+\rho K_s .\label{an04codv1:field}
\end{align}
BRW proves $Y>0$ from the positive remote entrance. On $s\ge0$ its
\texttt{budget} gives
$\max(Y,\rho K_s)\le B_0e^{Ls}(1+s)^{-2/3}$.
Here $S_x\le K_s$ by the nonnegative Gaussian overlap integral.
We also need two consequences within BRW's proof. To keep their
mass and time ranges explicit we give their derivations below.

\paragraph{Transfer sources.}
P \texttt{pa:exactflow}, \texttt{pa:earlybounds} and
\texttt{pa:targetcomparison} concern the exact physical Cartesian
dynamics and the small-cap comparison under its early alignment
guard. CRA \texttt{an04crav1:outer} supplies the target-only compact
Q2 law at the critical nominal cap; PROFILE
\texttt{inc:AN02:profile:v1:angle} supplies the target-only Gaussian
angular field on Q3. Their separate domains are checked in the
trained-flow proof. The finite-time limit uses BRW's delay and
extraction, with no weak-passage inputs.

\section{Two tail estimates on the logistic law}
\label{an04codv1:tail}
For Q1, $f=\pi\{\tilde\alpha\ge\alpha\}\in(0,1/2)$ and
$|\dd f|=\dd\pi$ on that sector. At $s=0$ HBD physical and
$m=2\cos^2\alpha$ give $x_++y_+\le C(1+f^{-1/3})$:
$f\asymp(\pi/2-\alpha)^3$ near the vertical axis, and each
limiting angle is bounded there by $C/\cos\alpha$. Q4 is bounded
using a single compact interior representative. Since $M\ge1/2$,
order implies $S_x\ge MfK(\rho x)$. For positive parts
\[
 2D^+\binom{x_+}{y_+}\le
 \lambda\begin{pmatrix}-f/2&1\\1&0\end{pmatrix}\binom{x_+}{y_+}
       +\rho\phi(0)\binom11.
\]
At a positive coordinate use \eqref{an04codv1:field}, $Y\ge0$,
$K(t)\ge t$, $K(t)\le\phi(0)+t_+$. At a negative coordinate its
positive part has zero derivative and the corresponding right side
is nonnegative; at equality with zero it bounds the positive-part
derivative as well. For Q4 take $f=1/2$. Multiply by the positive
left eigenvector $(w,1)$ of the matrix, $w=(\sqrt{4+(f/2)^2}-f/2)/2$.
It stays bounded away from zero and $Lw\le L-\lambda f/16$
by $\sqrt{4+d^2}\le2+d^2/4$ for $0\le d\le1$. Integration gives
\begin{equation}\label{an04codv1:ceiling}
 x_+(s)+y_+(s)\le C(1+f^{-1/3})e^{Ls-bfs},\qquad b=\lambda/16 .
\end{equation}
In particular $X_m\le\int x_+\dd\mu\le C e^{Ls}$; the sharper
BRW moment budget follows by integrating the damped ceiling.

For the lower estimate, HBD outer at half mass gives initial
$\min(x,y)\ge4^{-\omega}(c^\omega\tan\alpha-J_*)$ for high Q1
ranks. Thus $\pi\{\min(x(0),y(0))\ge R\}\ge c_1(1+R)^{-3}$
for large $R$, uniformly in the law. For a horizon $T$, put
$g=e^{Ls}\sqrt{M(0)/M(s)}$, $b_0>\sqrt2 B_0$,
$R_T=d+3b_0[(1+T)^{1/3}-1]$ for fixed sufficiently large $d$,
and $w_T=g(R_T-3b_0[(1+s)^{1/3}-1])$. On $0\le s\le T$,
$w_T\ge dg\ge de^{Ls}/\sqrt2$. Labels initially above $R_T$
in both coordinates satisfy both lower faces. On $x=w_T,y\ge w_T$
use $S_x\le K_s$ and $\rho\phi(\rho x)+\lambda yP_x\ge\rho K(\rho x)$.
On $y=w_T,x\ge w_T$ use monotonicity of $K$.
Both doubled velocities exceed $2w_T'$ since
$2g'/g=\lambda-(1-M)$, $\rho K(\rho w_T)>\lambda w_T$, and
$b_0g>B_0e^{Ls}$. This includes the corner and the initial faces.
At time $T$ the actual mass of this set is at least
$M(T)c_1(1+R_T)^{-3}\ge c_2/(1+T)$. Multiplying by the lower
input gives BRW's protected first input moment at least
$c_3D(T)$ for every $T\ge0$. The set varies with $T$.

\section{Core faces, mean output and exposure}
\paragraph{Upper core faces on every finite horizon.}
HBD gives $x_a(0),y_a(0)\le a$ at half mass. For $d=a+\epsilon$
the entry into $x_a,y_a\le d$ is strict. All donors are above the
bottom path, hence $X_m\ge M x_a$, $Y\ge M y_a$, and
$S_{x_a}=M\phi(\rho x_a)+\rho P_{x_a}X_m$. In particular
\[
 2x_a'=-(1-M)x_a+\rho(1-M)\phi(\rho x_a)
                  +\lambda P_{x_a}(y_a-X_m).
\]
On $x_a=d,y_a\le d$ this is at most $(1-M)[\rho K(\rho d)-d]<0$.
On $y_a=d,x_a\le d$, $2y_a'\le\rho K(\rho d)-d<0$.
Both inequalities hold at the corner. They are strict because
$M(s)<1$ for every finite $s$ and $r-\rho K(\rho r)$ is increasing
with zero at $a$. This excludes exit on any finite continuation;
send $\epsilon$ down to zero. Thus $x_a,y_a\le a$ on all $s\ge0$.

\paragraph{A lower bound on the true output mean.}
On the protected set $K(\rho x)\ge\rho x$, and elsewhere $K\ge0$.
Consequently $\rho K_s\ge\lambda c_3 D(s)$ for $s\ge0$.
For all $s\ge1$ it suffices to integrate the inequality
$2Y'\ge-Y+\lambda c_3D$ over the last unit of time.
To use one explicit constant through the initial unit, normalize
$c_3$ downwards once as follows. On the preceding $-1\le s\le0$,
HBD outer in its nominal clock $\tau=\log4+s$ puts a fixed positive
normalized mass $p$ at $x\ge1$: choose a compact Q1 band with
$c^\omega\tan\alpha>J_*+4^\omega$. There $M\ge1/(1+e)$.
Take $c_3$ to be the smaller of BRW's protected constant and
$p/(1+e)$. Then $\rho K_s\ge\lambda c_3 e^{Ls}/(1+s_+)$
on $s\ge-1$. In either interval $(1-2M)Y\ge-Y$ because $Y>0$.
Integrate from $s-1$, using $1+r_+\le1+s$ for $s-1\le r\le s$.
This proves, with the unmodified BRW constant also valid for $s\ge1$,
\begin{equation}\label{an04codv1:mean}
 Y(s)\ge k_Y D(s),\qquad
 k_Y=\frac{\lambda c_3}{1+2L}\left(1-e^{-(L+1/2)}\right)>0 .
\end{equation}

\paragraph{Core output.}
The core equation gives $2y_a'+(1-M)y_a\le a-Y$.
The exact integration kernel from $r$ to $s$ is
$\sqrt{M(r)/M(s)}\in[1/\sqrt2,1]$ on $0\le r\le s$ because
$(\log M)'=1-M$. Thus for $s\ge1$,
\[
 y_a(s)\le a+as/2-b_yD(s),\qquad b_y=k_Ye^{-L}/(2\sqrt2),
\]
where the negative contribution uses \eqref{an04codv1:mean} on
$[s-1,s]$. The polynomial term is eventually at most $b_yD(s)/2$.
Taking $k_y=b_y/2$ gives $y_a\le-k_yD$ for all sufficiently large $s$
and $y_a\le C-k_yD$ globally. In the opposite direction the same
kernel, BRW's upper bound on $Y$ and $\rho K\ge0$ give
$(y_a)_-\le C e^{Ls}$, using the uniform finite lower entry bound.

\paragraph{A lower input face.}
At $x_a<0$, the first two terms of its doubled equation are
nonnegative. The bounds just established for $(y_a)_-$ and $(X_m)_+$
therefore give $2x_a'\ge-Ce^{Ls}\Phi(\rho x_a)$.
Put $\ell(s)=-H\sqrt{1+s}$. Choose $H$ so the entry is strict
($x_a(0)$ has a uniform lower bound), $\rho^2H^2/2-L\ge1$ and
$Ce^{-\rho^2H^2/2}<H$. The Gaussian bound
$\Phi(-u)\le e^{-u^2/2}$ for $u\ge0$ follows by the exponential
moment inequality. On the sole lower face $x_a=\ell(s)$,
\[
 2x_a'\ge-Ce^{-\rho^2H^2/2}e^{-s}>
 -H/\sqrt{1+s}=2\ell'(s).
\]
There is no upper or output face in this scalar comparison.
Order gives $x\ge x_a>\ell(s)$ for all ranks. Hence for
$A_\zeta=\int(\zeta-x)_+\dd\mu$,
\begin{equation}\label{an04codv1:exposure}
 A_\zeta\le \zeta+H\sqrt{1+s},\qquad
 A_\zeta^2/2-\zeta A_\zeta\le C(1+s).
\end{equation}

\section{The active non-core upper bound}
\paragraph{First restriction: output at an active input.}
Set $h(s)=H_1e^{Ls}(b_1+s)^{-2/3}$, $b_1\ge\max(1,4/(3L))$.
Then $h'\ge Lh/2$ and $h$ is comparable to
$H_1 e^{Ls}(1+s)^{-2/3}$ with fixed constants.
We exclude the quadrant $x<\zeta,\ y>h(s)$, including the possibility
of entering it from an initially inactive label.
At entry, $x(0)\le\zeta$ on Q1 forces $f\ge f_*>0$: HBD outer
puts all sufficiently high Q1 labels strictly above $\zeta$.
By \eqref{an04codv1:ceiling} the remaining labels have $y(0)\le C_*$
there; use the core bound on Q4. Choose $h(0)>C_*+a+1$.
Enlarge $H_1$ so that at every $s\ge0$
\[
 \lambda\Phi(\rho\zeta) h(s)>\zeta+B_0e^{Ls}(1+s)^{-2/3},
 \qquad 2h'(s)>\rho K(\rho\zeta).
\]
Such a choice exists since $h$ is increasing and its ratio to the
budget envelope is bounded below by $H_1 b_1^{-2/3}$.
At the vertical boundary $x=\zeta,\ y\ge h$ one has
$2x'\ge-(1-M)\zeta-\rho K_s+\lambda\Phi(\rho\zeta)h>0$.
At the horizontal boundary $y=h,\ x\le\zeta$ one has
$2y'\le\rho K(\rho\zeta)<2h'$ using $Y>0$.
Both directions exclude the forbidden quadrant at the corner.
Continuity and these strict faces prove $y_+\le h(s)$ for every
currently active label.

\paragraph{Second restriction: mass capable of positive output.}
For large $s$ let $F_s=D_0\log(1+s)/(1+s)$ with $bD_0/2>3$,
where $b=\lambda/16$ in \eqref{an04codv1:ceiling}.
Consider the same label on the entire interval $s-1\le r\le s$;
its normalized rank $f$ stays fixed. If $f\ge F_s$, then
$r\ge(1+s)/2$ for $s\ge3$ and
\[
 \sup_{[s-1,s]}(x_++y_+)
 \le C(1+F_s^{-1/3}) e^{Ls}(1+s)^{-bD_0/2}=o(D(s))
\]
uniformly. Throughout that interval $Y(r)\ge k_Y e^{-L}D(s)$.
For all sufficiently large $s$, the preceding ceiling therefore
implies $\rho K(\rho x)\le\rho\phi(0)+\lambda x_+
\le k_Y e^{-L}D(s)/2$. While $y\ge0$,
\[
 y'\le-\sigma D(s),\qquad \sigma=k_Ye^{-L}/4.
\]
At entry $y_+(s-1)=o(D(s))<\sigma D(s)$, so a positive label
must cross zero before $s$; one initially at or below zero cannot
cross upwards because the same strict inequality holds at zero.
It follows that $y(s)>0$ on Q1 requires $f<F_s$.
This is an upper bound on positive-output ranks, whether active
or inactive. Their actual Q1 mass is at most $MF_s\le F_s$.

\paragraph{Combine the restrictions.}
In the active Q1 integral, the exposure of one label is at most
$\zeta+H\sqrt{1+s}$ by the core lower face, its positive output
is at most $h(s)$ by the quadrant argument, and its positive-output
mass is at most $F_s$. Thus
\begin{equation}\label{an04codv1:positive}
 0\le\int_{\mathrm{Q1},\,x<\zeta} y_+(\zeta-x)\dd\mu
 \le C e^{Ls}(1+s)^{-7/6}\log(1+s)=o(D(s)).
\end{equation}
No estimate of a single outlier is used in place of this integral.
For bounded $s$ it is uniformly bounded by the exposure and quadrant
bounds. Now split the full output term:
\[
 \int y(\zeta-x)_+\dd\mu
 =\tfrac M2 y_a(\zeta-x_a)
 +\int_{\mathrm{Q1},\,x<\zeta} y(\zeta-x)\dd\mu+0 .
\]
The last zero includes labels at the gate. For large $s$ the first
term is at most $-(\zeta-a)k_yD(s)/4$ since $y_a\le-k_yD$,
$M/2\ge1/4$, and $x_a\le a$. Combining \eqref{an04codv1:exposure},
\eqref{an04codv1:positive} and the exact leading formula
$\mathcal E_\zeta=A_\zeta^2/2-\zeta A_\zeta+\int y(\zeta-x)_+\dd\mu$
proves Theorem~\ref{an04codv1:main}, for example with
$k_E=(\zeta-a)k_y/8$ after increasing $C$ on a bounded interval.

\section{Transfer of the probe at each fixed resident time}
\begin{theorem}[Original trained flow]\label{an04codv1:trained}
For each fixed finite $s\ge0$,
\[
 \limsup_{q\downarrow0}\frac{E_\xi(t_h+s)-1/2}{q^2}
       \le C-k_E D(s).
\]
More precisely, along each BRW extracted subsequence the scaled
error converges to the probe functional of that extracted law
at $s$. For every $K_0>0$ and fixed weak response level $\eta>0$,
there exist a finite $s_{K_0}$ and a $q_0>0$ such that for every
$0<q<q_0$,
\[
 t_h+s_{K_0}<t_\eta,\qquad
 \frac{E_\xi(t_h+s_{K_0})-1/2}{q^2}<-K_0 .
\]
Here $t_\eta$ is the first post-half-mass hit of the BRW weak-response
threshold, possibly infinite. In particular the infimum on
$[t_h,t_\eta)$ tends to $-\infty$ in this scaled convention.
\end{theorem}
\begin{proof}
In Cartesian notation $f(\xi)=\int W(U^T\xi)_+\dd l$,
$\dd l=\dd\alpha/(2\pi)$, $\xi=(\sin(q\zeta),-\cos(q\zeta))$.
Fix $s$ throughout this transfer. BRW's finite-horizon physical
bound gives $|U|,|W|\le C_s$ and $|U_2|,|W_2|\le C_s q$ at
observation and on any fixed nominal-clock interval ending there
and starting at a fixed positive cap. Thus the individual-label
integrands of $H_q=f_1(\xi)/q$ and $V_q=f_2(\xi)/q^2$ are
both bounded by $C_s$ on the entire circle, since $|U^T\xi|\le C_s q$.
An initial-angle strip of length $\delta$ costs $O_s(\delta)$
in each scaled observable.

On compact right interiors, HBD--BRW give bounded $x,y$ and
convergence of the joint label-path law, including $m$, along extraction.
The exact integrands are $m\cos(qy)[\sin(q(\zeta-x))]_+/q$ and
$m[\sin(qy)/q][\sin(q(\zeta-x))]_+/q$.
Their limits are $m(\zeta-x)_+$ and $my(\zeta-x)_+$ on these
compacts. Positive part is continuous even for a gate atom.
These formulas transfer at $t_h+s$ by HBD clock and compact time
convergence. Endpoint-strip control also applies to the limiting
integrals by taking limits on compact strips first.

\emph{Left ancestry at the required output precision.}
A mass estimate $o(q)$ alone would not justify discarding $V_q$.
Fix a sufficiently small nominal cap $A_0$, at time
$T_0=\log(A_0/S_0)$. P \texttt{targetcomparison} applies there
with angles within $O(A_0)$ of the target-only angle and relative
radial weights within factors $e^{O(A_0)}$. The time from $T_0$ to
$t_h+s$ is bounded for fixed $s,A_0$ by HBD clock. On this interval
the generic radial bound $|(\log m)'|\le(1+2q^2)(1+\int m\dd l)$
and the physical envelope give a bounded radial multiplier.
On a compact interior Q2 arc, CRA \texttt{outer} applies with
$\kappa=1/\omega$ and $\kappa\lambda/2=\lambda/(1-\lambda)<1$.
At $T_0$ the target angle tends to $\pi/2$ and its radial weight
is comparable to $S_0e^{\lambda T_0}\asymp q^2$.
Take $A_0$ small before $q$. Both trained second components are
then at least $k_I q$ at this cap; $|U|,|W|\le C_{s,I}q$
up to observation.

To check the signs up to observation, the Cartesian field is
$U'=R^TW$, $W'=RU$, $R=A_\theta-C_f$.
The physical bounds above give $|f_1(D)|\le C_s(|D_1|+q|D_2|)$
and $|f_2(D)|\le C_s(q|D_1|+q^2|D_2|)$ for any data vector $D$.
Since $\mathbb E|D_1D_2|\le Cq$ under the two Gaussian clusters,
$|R_{12}|+|R_{21}|\le C_s q$ and $R_{22}\ge-C_s q^2$
(the target diagonal is nonnegative). While $U_2,W_2\ge0$ these
estimates give $U_2',W_2'\ge-C_{s,I}q^2$. Integration from
their strict entry excludes a first contact of either component
with $k_I q/2$ over this fixed interval, including simultaneous
contact. Consequently $U^T\xi<0$ on the Q2 compact at observation.

On a compact interior Q3 arc both target-only directions have
negative projections on both cluster means, separated from zero
at initialization. They retain half of this separation until $T_0$:
on that neighborhood the exact target angular drift is
$q[-\sin\theta K(\cos\theta/q)+\rho\cos\theta
K(\rho\sin\theta/q)]/2$ (PROFILE \texttt{angle}), exponentially
small by $K(-v)\le\phi(v)$ for $v\ge0$. Its integral over
$O(\log(1/q))$ cannot reach either enlarged arc endpoint.
The target log radial rate $2\mathbb E[(a^TD)_+^2]$ is exponentially
small there by the Gaussian tail and Cauchy--Schwarz. The trained
radial weight at $T_0$ is therefore $O_I(S_0)$ by P, and remains
$O_{s,I}(S_0)$ at observation. Its output divided by $q^2$
vanishes since $S_0/q^2\to0$. This treats both left interiors.
The axis labels have zero initial measure; take $q\downarrow0$ on
the interior compacts and then send all strip lengths to zero.

It follows that $(H_q,V_q)\to(A_\zeta,\int y(\zeta-x)_+\dd\mu)$
on the extracted subsequence. The exact energy identity is
\[
 (E_\xi-1/2)/q^2=H_q^2/2-[\sin(q\zeta)/q]H_q
                 +\cos(q\zeta)V_q+q^2 V_q^2/2 .
\]
This proves convergence on extraction at each fixed $s$. Every
subsequence has a further extraction and the limiting upper bound
is uniform in the law, proving the full limsup. Choose $s_{K_0}$
so that $C-k_ED(s_{K_0})<-K_0-1$, and only then make $q$ small.
BRW delay gives $\sup_{0\le r\le s_{K_0}}|\mathsf m_2(t_h+r)|\to0$.
This proves the strict pre-weak witness and the assertion about
the infimum without requiring attainment at the right endpoint.
\end{proof}

\paragraph{Scope and constants.}
Uniformity over laws uses the fixed-$c,\rho,\zeta$ entry, tail and
budget constants from HBD and BRW. In particular $c_1,c_2,c_3$,
the input-face constant, and the onset of the rank cutoff may
depend on $c$; no common depth rate as $c\downarrow0$ is claimed.
The mean $Y$ supplies the negative output drive on the active core,
as in MPB's bridge mechanism; here the growing mean has an actual
protected-tail source on the pre-weak time scale. No bound is
asserted at $s=t_\eta-t_h$. The post-learning swing-by minimum is
a different comparison.
\end{document}
